\noindent
{\Large \textbf{\section{Результаты работы}}} \par

Перевод вещественного числа $-123.45$ в формат bf16 и обратно, с разбором результата по битам представлен в Листинге~\ref{lst:result-bf16}. Так как мантисса bf16 ограничена 7 битами, после округления восстановленное число немного отличается от исходного ($-123.5$ вместо $-123.45$).

\begin{lstlisting}[language=,caption=Перевод вещественного числа -123.45 в bf16 туда и обратно,label=lst:result-bf16]
Выберите направление перевода:
  1. десятичное -> двоичное
  2. двоичное -> десятичное
  0. выход
> 1
Тип числа:
  1. целое
  2. вещественное (bf16)
> 2
Введите вещественное десятичное число: -123.45
Двоичный результат (bf16, 16 бит): 1100001011110111
знак = 1 (отрицательное)
экспонента = 10000101 -> смещённая 133, истинная 6 (смещённая - 127)
мантисса = 1110111 -> дробная часть 0.9296875

Выберите направление перевода:
  1. десятичное -> двоичное
  2. двоичное -> десятичное
  0. выход
> 2
Тип числа:
  1. целое
  2. вещественное (bf16, 16 бит)
> 2
Введите 16 бит числа в формате bf16: 1100001011110111
знак = 1 (отрицательное)
экспонента = 10000101 -> смещённая 133, истинная 6 (смещённая - 127)
мантисса = 1110111 -> дробная часть 0.9296875
Десятичный результат: -123.5
\end{lstlisting}

\newpage
